\documentclass[11pt]{article}
\usepackage[margin=0.8in]{geometry}
\usepackage{amsmath, amssymb}
\usepackage{booktabs}
\usepackage{array}
\usepackage{enumitem}
\usepackage{fancyhdr}
\usepackage{parskip}
\usepackage{bm}

% Header and Footer
\pagestyle{fancy}
\setlength{\headheight}{13.6pt}
\fancyhf{}
\fancyhead[C]{\textbf{Chapter 3: Propagation of Error Summary Notes}}
\fancyfoot[C]{CEE 305 - Spring 2026}
\renewcommand{\headrulewidth}{1pt}
\renewcommand{\footrulewidth}{0.5pt}

\begin{document}

% ==========================================
\section*{Introduction to Error Propagation}
% ==========================================
\textbf{Key Concept:} Any measuring procedure contains error. Measured values generally differ from the true values being measured, and these errors produce errors in calculated values (such as the mean).

\begin{quote}
\textbf{Definition:} When error in measurement produces error in calculated values, we say that \textbf{error is propagated} from the measurements to the calculated value.
\end{quote}

% ==========================================
\section*{Section 3.1: Measurement Error}
% ==========================================
\textbf{Basic Definition:} The \textbf{error} in a measured value is the difference between a measured value and the true value.

\textbf{Components of Error:}
\begin{itemize}[nosep]
    \item \textbf{Systematic error (bias)} – The part of the error that is the same for every measurement.
    \item \textbf{Random error} – Error that varies from measurement to measurement and averages to zero in the long run.
\end{itemize}

\textbf{Mathematical Representation:}
\[
\text{Measured value} = \text{true value} + \text{bias} + \text{random error}
\]

\textbf{Key Properties of Measuring Processes:}
\begin{itemize}[nosep]
    \item \textbf{Accuracy} – Determined by bias (smaller bias = more accurate).
    \item \textbf{Unbiased process} – When bias = 0.
    \item \textbf{Precision} – The degree to which repeated measurements tend to agree (higher precision = measurements cluster closely).
\end{itemize}

\subsection*{Statistical Characterization of Measurements}
A measured value is a random variable with:
\begin{itemize}[nosep]
    \item \textbf{Mean (\(\mu\))} – The average value over many measurements.
    \item \textbf{Standard deviation (\(\sigma\))} – Measures the spread of measurements.
\end{itemize}

\textbf{Bias and Uncertainty:}
\[
\text{Bias} = \mu - \text{true value},\qquad \text{Uncertainty} = \sigma
\]
\begin{itemize}[nosep]
    \item Smaller bias \(\rightarrow\) more accurate process.
    \item Smaller uncertainty \(\rightarrow\) more precise process.
\end{itemize}

\subsection*{Estimating Bias and Uncertainty}
Given independent measurements \(X_1, X_2, \ldots, X_n\) of the same quantity:
\begin{itemize}[nosep]
    \item \textbf{Estimate uncertainty} using sample standard deviation:
    \[
    s = \sqrt{\frac{1}{n-1}\sum_{i=1}^n (X_i - \bar{X})^2}
    \]
    \item \textbf{Estimate bias} (if true value is known):
    \[
    \text{Bias} \approx \bar{X} - \text{true value}
    \]
    \item \textbf{Note:} If true value is unknown, bias cannot be estimated from repeated measurements alone.
\end{itemize}

% ==========================================
\section*{Section 3.2: Linear Combinations of Measurements}
% ==========================================
\textbf{Single Measurement Scaled by Constant:}
If \(X\) is a measurement and \(c\) is a constant:
\[
\sigma_{cX} = |c| \sigma_X
\]

\textbf{Linear Combination of Independent Measurements:}
If \(X_1, \dots, X_n\) are independent measurements and \(c_1, \dots, c_n\) are constants:
\[
\sigma_{c_1X_1 + \cdots + c_nX_n} = \sqrt{c_1^2\sigma_{X_1}^2 + \cdots + c_n^2\sigma_{X_n}^2}
\]

\subsection*{Repeated Measurements}
If \(X_1, \dots, X_n\) are \(n\) independent measurements, each with mean \(\mu\) and standard deviation \(\sigma\):
\[
\mu_{\bar{X}} = \mu,\qquad \sigma_{\bar{X}} = \frac{\sigma}{\sqrt{n}}
\]
\textbf{Interpretation:} The average of \(n\) measurements is more precise than a single measurement by a factor of \(\sqrt{n}\).

\subsection*{Weighted Average for Measurements with Different Uncertainties}
If \(X\) and \(Y\) are independent measurements of the same quantity, with uncertainties \(\sigma_X\) and \(\sigma_Y\):

The weighted average with the smallest uncertainty is:
\[
c_{\text{best}}X + (1 - c_{\text{best}})Y
\]
where
\[
c_{\text{best}} = \frac{\sigma_Y^2}{\sigma_X^2 + \sigma_Y^2},\qquad 1 - c_{\text{best}} = \frac{\sigma_X^2}{\sigma_X^2 + \sigma_Y^2}
\]

\subsection*{Linear Combinations of Dependent Measurements}
If \(X_1, \ldots, X_n\) are measurements (possibly dependent) and \(c_1, \ldots, c_n\) are constants:
\[
\sigma_{c_1X_1 + \cdots + c_nX_n} \leq |c_1|\sigma_{X_1} + \cdots + |c_n|\sigma_{X_n}
\]
This provides a \textbf{conservative estimate} (upper bound) when dependencies are unknown.

% ==========================================
\section*{Section 3.3: Uncertainties for Functions of One Measurement}
% ==========================================
\textbf{Propagation of Error Formula (One Variable):}
If \(X\) is a measurement with small uncertainty \(\sigma_X\), and \(U = U(X)\) is a function of \(X\):
\[
\sigma_U \approx \left| \frac{dU}{dX} \right| \sigma_X
\]
In practice, evaluate \(dU/dX\) at the observed measurement \(X\).

\subsection*{Relative Uncertainty}
\textbf{Definition:} If \(U\) is a measurement with true value \(\mu_U\) and uncertainty \(\sigma_U\):
\[
\text{Relative uncertainty} = \frac{\sigma_U}{\mu_U}
\]
\textbf{Properties:}
\begin{itemize}[nosep]
    \item Unitless quantity (often expressed as a percentage).
    \item In practice, \(\mu_U\) is unknown.
    \item If bias is negligible, estimate relative uncertainty as \(\sigma_U / U\).
\end{itemize}

\textbf{Two Methods to Find Relative Uncertainty:}
\begin{enumerate}[nosep]
    \item Compute \(\sigma_U\) using \(\sigma_U \approx |dU/dX|\sigma_X\), then divide by \(U\).
    \item Compute \(\ln U\) and find \(\sigma_{\ln U}\) (which equals \(\sigma_U/U\)).
\end{enumerate}

% ==========================================
\section*{Section 3.4: Uncertainties for Functions of Several Measurements}
% ==========================================
\textbf{Propagation of Error Formula (Multiple Independent Variables):}
If \(X_1, \ldots, X_n\) are \textbf{independent} measurements with small uncertainties \(\sigma_{X_i}\), and \(U = U(X_1, \ldots, X_n)\):
\[
\sigma_U \approx \sqrt{\left(\frac{\partial U}{\partial X_1}\right)^2\sigma_{X_1}^2 + \left(\frac{\partial U}{\partial X_2}\right)^2\sigma_{X_2}^2 + \cdots + \left(\frac{\partial U}{\partial X_n}\right)^2\sigma_{X_n}^2}
\]
In practice, evaluate partial derivatives at \((X_1, \ldots, X_n)\).

\textbf{Conservative Estimate for Dependent Measurements:}
If \(X_1, \ldots, X_n\) are \textbf{dependent} measurements with small uncertainties:
\[
\sigma_U \leq \left|\frac{\partial U}{\partial X_1}\right|\sigma_{X_1} + \left|\frac{\partial U}{\partial X_2}\right|\sigma_{X_2} + \cdots + \left|\frac{\partial U}{\partial X_n}\right|\sigma_{X_n}
\]
\textit{Note:} This inequality holds in almost all practical situations (may fail if second partial derivatives are quite large).

% ==========================================
\section*{Common Applications and Examples}
% ==========================================
\textbf{Example: Volume of a Sphere} – Given radius \(R\) with uncertainty \(\sigma_R\):
\[
V = \frac{4}{3}\pi R^3,\quad \frac{dV}{dR} = 4\pi R^2,\quad \sigma_V = 4\pi R^2\sigma_R
\]

\textbf{Example: Area of a Rectangle} – Given sides \(X\) and \(Y\) with uncertainties \(\sigma_X\) and \(\sigma_Y\):
\[
A = XY,\quad \frac{\partial A}{\partial X} = Y,\quad \frac{\partial A}{\partial Y} = X,\quad \sigma_A = \sqrt{Y^2\sigma_X^2 + X^2\sigma_Y^2}
\]

% ==========================================
\section*{Summary Table of Error Propagation Formulas}
% ==========================================
{\small
\begin{center}
\begin{tabular}{l c}
\toprule
\textbf{Situation} & \textbf{Formula} \\
\midrule
Single measurement scaled by constant & \(\sigma_{cX} = |c|\sigma_X\) \\
Linear combination (independent) & \(\sigma_{\sum c_iX_i} = \sqrt{\sum c_i^2\sigma_{X_i}^2}\) \\
Linear combination (dependent) & \(\sigma_{\sum c_iX_i} \leq \sum |c_i|\sigma_{X_i}\) \\
Sample mean (\(n\) measurements) & \(\sigma_{\bar{X}} = \sigma/\sqrt{n}\) \\
Function of one variable & \(\sigma_U \approx |dU/dX|\sigma_X\) \\
Function of multiple independent variables & \(\sigma_U \approx \sqrt{\sum \left(\frac{\partial U}{\partial X_i}\right)^2\sigma_{X_i}^2}\) \\
Function of multiple dependent variables & \(\sigma_U \leq \sum \left|\frac{\partial U}{\partial X_i}\right|\sigma_{X_i}\) \\
\bottomrule
\end{tabular}
\end{center}
}

% ==========================================
\section*{Key Terminology}
% ==========================================
\begin{itemize}[nosep]
    \item \textbf{Error:} Difference between measured value and true value.
    \item \textbf{Bias:} Systematic component of error (constant across measurements).
    \item \textbf{Random error:} Variable component of error (averages to zero).
    \item \textbf{Accuracy:} How close measurements are to true value (related to bias).
    \item \textbf{Precision:} How close repeated measurements are to each other (related to uncertainty).
    \item \textbf{Uncertainty:} Standard deviation of measurements (\(\sigma\)).
    \item \textbf{Propagation of error:} How measurement errors affect calculated quantities.
    \item \textbf{Relative uncertainty:} Uncertainty expressed as a fraction of the value (\(\sigma_U/\mu_U\)).
\end{itemize}

% ==========================================
\section*{Important Notes}
% ==========================================
\begin{itemize}[nosep]
    \item All propagation formulas are approximations valid when uncertainties are small.
    \item Partial derivatives should be evaluated at the measured values.
    \item For dependent measurements, the sum of absolute values gives a conservative (upper bound) estimate.
    \item The sample mean reduces uncertainty by a factor of \(\sqrt{n}\) compared to single measurements.
    \item Weighted averages can minimize uncertainty when combining measurements with different precisions.
\end{itemize}

\end{document}